\documentclass[10pt,twocolumn]{article}
\usepackage[margin=0.8in,top=0.75in,bottom=0.75in]{geometry}
\usepackage{graphicx}
\usepackage{booktabs}
\usepackage{amsmath}
\usepackage{hyperref}
\usepackage{microtype}
\usepackage{parskip}
\setlength{\parskip}{4pt}

\title{\textbf{Bayesian Networks for Heart Disease Diagnosis:\\
Structure Learning, Calibrated Uncertainty, and Interpretation}}
\author{CS 179 Project Report\\University of California, Irvine}
\date{June 2026}

\begin{document}
\maketitle

\section{Introduction}
Heart disease is the leading cause of death in the United States, yet its
diagnosis involves integrating uncertain evidence from many heterogeneous
clinical variables. A \textit{probabilistic graphical model} is a natural
fit: a \textbf{Bayesian network} (BN) over the observed variables explicitly
represents conditional independence structure, supports exact inference over any
query variable, and---crucially---outputs calibrated posterior probabilities
rather than mere point predictions.

We apply BNs to the UCI Cleveland Heart Disease dataset (303 patients,
13 clinical features, binary target). Our primary questions are: (1) Does
a hand-crafted expert structure or a data-learned structure yield better
predictive performance? (2) Are the predicted probabilities \textit{well
calibrated}? (3) Which variables are most informative for diagnosis?

\section{Data and Methods}

\subsection{Dataset}
The Cleveland subset of the UCI Heart Disease dataset contains 303 patients
(68\% male, 55\% with diagnosed heart disease). The 13 features span demographics
(\texttt{age}, \texttt{sex}), risk factors (\texttt{chol}, \texttt{trestbps},
\texttt{fbs}), and clinical test results (\texttt{cp}, \texttt{restecg},
\texttt{thalach}, \texttt{exang}, \texttt{oldpeak}, \texttt{slope}, \texttt{ca},
\texttt{thal}).

\subsection{Discretisation}
Bayesian networks require discrete state spaces. Five continuous variables were
binned into clinically interpretable categories following established guidelines
(AHA blood-pressure tiers, NCEP cholesterol tiers, etc.):
\texttt{age} $\to$ 4 bins, \texttt{trestbps} $\to$ 3 bins, \texttt{chol} $\to$ 3 bins,
\texttt{thalach} $\to$ 3 bins, \texttt{oldpeak} $\to$ 3 bins.
The 8 remaining variables are already categorical or binary.

An 80/20 stratified split gives a training set of 242 and test set of 61 patients.

\subsection{Model 1: Expert DAG}
We encode prior clinical knowledge as a directed acyclic graph (DAG) with 17
edges. Demographics (\texttt{age}, \texttt{sex}) are root causes of
cholesterol, resting blood pressure, and maximum heart rate. The ECG result
(\texttt{restecg}) is a child of \texttt{trestbps}. The exercise pathway chains
\texttt{thalach} $\to$ \texttt{exang} $\to$ \texttt{target}. The ST-segment chain
is $\texttt{slope} \to \texttt{oldpeak} \to \texttt{target}$. Eight variables
have direct edges into \texttt{target}.

\subsection{Model 2: Learned DAG (Hill-Climbing / BIC)}
We use greedy hill-climbing \cite{pgmpy} with the Bayesian Information
Criterion as the scoring function, constraining maximum in-degree to 3 to keep
CPTs tractable. The search is run over the full set of 14 discretised variables
on the training fold.

\subsection{Parameter Learning}
Both structures are fitted by Maximum Likelihood with a BDeu Bayesian smoothing
prior (equivalent sample size = 5) to avoid zero-probability CPT entries on the
small dataset.

\subsection{Inference}
We use exact inference via Variable Elimination with the MinFill elimination
order. For each test patient we compute P(target $= 1 \mid$ all other 13
features observed).

\subsection{Evaluation Metrics}
\begin{itemize}\setlength{\itemsep}{0pt}
  \item \textbf{Accuracy} at the 0.5 threshold.
  \item \textbf{ROC-AUC} for ranking quality independent of threshold.
  \item \textbf{Brier score} $= \frac{1}{n}\sum(\hat{p}_i - y_i)^2$,
        which jointly measures calibration and resolution.
  \item \textbf{Calibration curve} (reliability diagram): empirical positive
        fraction vs.\ predicted probability in equal-width bins.
  \item \textbf{5-fold stratified cross-validation} for stable generalisation
        estimates.
\end{itemize}

\section{Results and Discussion}

\subsection{Learned Structure}
Figure~1 shows both DAGs. Hill-climbing recovers several edges present in the
expert graph (\texttt{ca}$\to$\texttt{target},
\texttt{cp}$\to$\texttt{target}, \texttt{thal}$\to$\texttt{target},
\texttt{oldpeak}$\to$\texttt{target}) and adds data-driven connections that
were not in the expert prior. The expert structure has higher BIC on
training data in some runs and lower in others depending on how aggressively
the search explores, illustrating the classic bias-variance tradeoff between
prior knowledge and data.

\subsection{Predictive Performance}

\begin{table}[h]
\centering\small
\caption{Test-set (n=61) and 5-fold CV metrics.}
\begin{tabular}{lrrr}
\toprule
Model & Accuracy & ROC-AUC & Brier \\
\midrule
Expert DAG        & 0.820 & 0.882 & 0.140 \\
Learned DAG (HC)  & 0.803 & 0.869 & 0.148 \\
\bottomrule
\end{tabular}
\end{table}

\noindent Both models achieve strong discrimination (AUC $> 0.86$), with
the expert DAG marginally outperforming the data-driven one. This is a
common finding on small datasets: domain knowledge acts as a useful
structural regulariser and prevents the search from overfitting spurious
dependencies.

\subsection{Calibration}
Figure~2 shows reliability diagrams for both models. The expert DAG stays
closer to the diagonal, yielding a lower Brier score and indicating that
its predicted probabilities are better calibrated. The learned model
tends to be slightly overconfident in the high-probability bins---a known
consequence of maximum-likelihood structure learning without adequate
regularisation.

Well-calibrated probabilities matter in a clinical context: a doctor
receiving ``80\% probability of heart disease'' should be able to trust
that approximately 80\% of similarly-scored patients do have the disease.
BNs with a BDeu prior approach this calibration naturally, in contrast to
a discriminative classifier that would require post-hoc Platt scaling.

\subsection{Learning Curves}
Figure~3 shows ROC-AUC and Brier score as training-set size grows from 20\%
to 100\% of the training data. AUC rises monotonically and stabilises
above 200 samples, while Brier score correspondingly decreases, confirming
that the dataset is large enough to fit the 17-edge expert model reliably.

\subsection{Interpretation and Markov Blanket}
The Markov blanket of \texttt{target} in the expert model consists of all
8 direct parents plus the children of \texttt{target} (none in our expert
DAG) and co-parents of those children. These 8 variables---\texttt{ca},
\texttt{cp}, \texttt{thal}, \texttt{oldpeak}, \texttt{exang},
\texttt{chol}, \texttt{trestbps}, \texttt{restecg}---match the clinical
literature: number of blocked vessels (\texttt{ca}), chest-pain type
(\texttt{cp}), thalassemia type (\texttt{thal}), and ST depression
(\texttt{oldpeak}) are widely accepted as the strongest predictors
of coronary artery disease.

The learned model recovers most of these high-weight edges, providing
data-driven validation of the expert structure. Edges that appear only in
the learned graph (e.g., \texttt{fbs}$\to$\texttt{target}) represent
relationships that are statistically present in this sample but lack strong
prior clinical justification.

\subsection{5-Fold Cross-Validation}
Over 5 folds, the expert DAG achieves AUC $= 0.874 \pm 0.038$ and the
learned DAG $0.862 \pm 0.051$. The learned model has higher variance across
folds, consistent with its dependence on the specific training sample.

\section{Conclusion}
We have shown that a Bayesian network with a clinically informed structure
can accurately predict heart disease (AUC $\approx 0.88$, accuracy $\approx 82\%$)
with well-calibrated posterior probabilities. The expert DAG outperforms
the fully data-driven structure on all metrics when training data is
limited---a predictable consequence of structure search overfitting to
242 samples. Nonetheless, the learned model independently recovers the
strongest clinical predictors, providing data-driven support for the
expert encoding. The calibration analysis demonstrates that BNs with
BDeu smoothing are near-perfectly reliable at intermediate confidence
levels and only slightly overconfident at high probabilities, making them
suitable for decision-support in clinical settings where uncertainty
quantification matters.

Future work could explore Chow-Liu trees for continuous variables (avoiding
discretisation), latent variable extensions to model unobserved comorbidities,
or comparison with variational inference in a Pyro BN to scale to larger
datasets.

\begin{thebibliography}{9}
\bibitem{pgmpy}
  A. Ankan, A. Panda, \textit{pgmpy: Probabilistic Graphical Models using
  Python}, SciPy 2015. \url{https://pgmpy.org}

\bibitem{uci}
  D. Janosi et al., \textit{Heart Disease Data Set}, UCI ML Repository, 1988.
  \url{https://archive.ics.uci.edu/ml/datasets/heart+disease}

\bibitem{hc}
  D. Chickering, \textit{Optimal Structure Identification With Greedy Search},
  JMLR 3, 507--554, 2002.

\bibitem{bdeu}
  W. Buntine, \textit{Theory Refinement on Bayesian Networks},
  UAI 1991.
\end{thebibliography}

\end{document}
